\section{Compressed surface covers and ordered attachment points}
\label{sec:surface-covers}

Report 24's second contribution concerns geometric compression rather than
chain cancellation.  Its dihedral-cover classifier is now maintained in
\code{surface\_cover.py}.  We also extend it with exact comparison of
components carrying an ordered list of points over a fixed basepoint.
This is a possible interface for part of the attachment bookkeeping in a
future hierarchy algorithm.  It currently receives a covering presentation;
it does not extract one from a knot exterior and supplies no knot verdict.

The general correspondence between connected covers and transitive
fundamental-group actions, with equivariant bijections representing cover
isomorphisms, is classical~\cite[Section 1.3]{hatcher-covering}.  The formulas
below specialize it to binary affine data.  They are not a replacement for
general normal-surface orbit algorithms: Agol--Hass--Thurston already give
polynomial algorithms for orbit counting and associated normal-surface
topology~\cite{aht-orbits}.  Our new marked signatures have the narrower
global dihedral input promise stated here.

\subsection{A checked presentation and three unmarked families}

Let $F$ be a connected compact surface with $b\ge1$ boundary circles and
free fundamental group of rank $r=1-\chi(F)$.  The input specifies its
orientability, genus (crosscap number when nonorientable), and $b$, together
with $r$ monodromies on a fibre of binary size $W$:
\[
 g_i(x)=\epsilon_i x+a_i\pmod W,\qquad \epsilon_i\in\{1,-1\}.
\]
Each assignment defines a cover because this fundamental group is free.
Peripheral words are generated from the canonical surface schema: the
last is the inverse of the product of handle commutators, or crosscap
squares, followed by the first $b-1$ boundary generators.  Arbitrary supplied
words are not accepted as a peripheral-system certificate.  Closed bases
are outside this interface.  The disk is included.

If all signs are positive, put $d=\gcd(W,a_1,\ldots,a_r)$ and $m=W/d$.
Otherwise fix $R(x)=-x+a_0$ and take the gcd of $W$, all positive shifts,
and all differences $a_i-a_0$ for negative generators.  In either case
the translation subgroup is generated by $T(x)=x+d$.  Products of two
reflections give the differences, and conjugating a translation by a
reflection negates its shift; Bezout supplies the converse containment.

Without reflections there are $d$ components of degree $m$, all isomorphic
over $F$.  With a reflection, components are the orbits of
$r_0\mapsto a_0-r_0$ on residues modulo $d$.  Let $f$ be the number of
solutions of $2r_0=a_0\pmod d$.  There are $(d-f)/2$ paired-residue
components of degree $2m$ and $f\le2$ fixed-residue components of degree $m$.
All paired components are equivariantly isomorphic: for two nonfixed
residues $r_0,s_0$, send
\[
 r_0+jd\mapsto s_0+jd,\qquad
 a_0-r_0+jd\mapsto a_0-s_0+jd.
\]
These maps commute with $T,R$, hence with every monodromy generator.
Thus at most three records, with binary multiplicities, suffice.

When two fixed residues occur, their covers are isomorphic exactly when
$m$ is odd.  An equivariant map between their transitive translation
orbits must be a translation by $A=d/2+dt$.  Commutation with $R$ requires
$2A=0\pmod W$, or $2t=-1\pmod m$.  This is solvable exactly for odd $m$.
The records retain distinct component labels even when their cover types
coincide.  Type equality never means that external attachments can be
discarded.  The three-type bound is sharp: for a genus-one base with two
boundaries and maps $x+d,-x,-x$, even $d\ge4,m\ge2$ give component genera
$m+1,m/2,m/2+1$ and boundary counts $2m,m+2,m$.

\subsection{Orientation and peripheral information are retained}

An affine sign is not the orientation character of a surface loop.
Attach the actual character $w_i\in\mathbb Z/2$ to every generator.
For the translation increments $u_j$, retain bits $v_j=w_j$ for positive
generators and $v_j=w_j+w_0$ for reflection differences.  An extended gcd
identity gives $(d,\tau)$ as a combination of $(W,0)$ and $(u_j,v_j)$.
The signed translation action is consistent precisely when
\[
 m\tau=0\pmod2,\qquad v_j=(u_j/d)\tau\pmod2\quad\hbox{for every }j.
\]
If a relation fails, subtracting multiples of $(d,\tau)$ produces
$(0,1)$, an orientation-reversing loop fixing every sheet.  Otherwise the
signed translation group is $(jd,j\tau)$, with no orientation-reversing
translation stabilizer.  A reflection stabilizes a fixed-residue component
at $r_0$ with orientation bit
$w_0+((2r_0-a_0)/d)\tau$.  Thus all paired components are orientable in
the consistent case, and this last bit decides each fixed component.
Pure translation covers use the same consistency test.  In particular,
an even-degree connected cyclic cover of a M\"obius band is an annulus;
an odd-degree one remains a M\"obius band.

For boundary monodromy $x+c$, put $q=\gcd(m,c/d)$.  There are $q$ lifts
of degree $m/q$ on a single-residue component, or $2q$ on a paired one.
A boundary reflection exchanges the two residues of a paired component,
giving $m$ lifts of degree two.  On a fixed residue $r_0$, its fixed
points solve $2y=(c-2r_0)/d\pmod m$.  If this has $\nu$ solutions, the
profile is $\nu$ degree-one and $(m-\nu)/2$ degree-two lifts, omitting
zero multiplicities.  Together with
\[
 \chi(C)=\deg(C)\chi(F),\qquad
 2-\#\partial C-\chi(C)=
 \begin{cases}2g(C),&C\text{ orientable},\\g(C),&C\text{ nonorientable},\end{cases}
\]
these profiles determine the topology.  The code checks the resulting
integer identities.  They describe the orientation interval bundle with
orientable total space; a different bundle character would be extra input.

The component of sheet $x$ is labeled by $x\bmod d$, together with its
reflection partner when distinct.  A boundary translation cycle has label
$x\bmod\gcd(W,c)$; a reflection cycle has label $\{x,c-x\pmod W\}$.
These give individual component and boundary-lift queries without
enumerating other sheets or lifts.

\subsection{An exact extension for ordered marked points}

Fix the normalized input presentation.  A marked component consists of
one of its components and an ordered list $(x_1,\ldots,x_q)$ of fibre
points in that component.  Repeated points are allowed.  Equivalence means
an isomorphism over the identity of $F$ taking each $x_i$ to its corresponding
point.  The implementation's component-selector sheet identifies the
component and is not an additional mark.  This definition does not include
unparameterized boundary-circle markings or arbitrary attaching maps.

\begin{proposition}[Rooted equivariant transport]
An equivariant map between connected components is uniquely determined by
the image $y$ of one point $x$.  For two translation components any $x,y$
are admissible and the map is $z\mapsto z-x+y$.  For two paired components,
any $x,y$ are admissible and the map is
\[
 f_{x,y}(z)=
 \begin{cases}
 z-x+y,&z\equiv x\pmod d,\\
 z+x-y,&z\equiv a_0-x\pmod d,
 \end{cases}\pmod W.
\]
For two fixed components it exists precisely when
$2(y-x)=0\pmod W$, and is then the translation $z\mapsto z-x+y$.
A paired and a fixed component cannot be isomorphic over the base.
\end{proposition}
\begin{proof}
Transitivity forces $f(gx)=gy$ for every monodromy element $g$.
Translations act regularly on a single residue class.  On a paired
component, the points have unique expressions $T^j x$ or $T^jRx$, for
$j\in\mathbb Z/m$; the two displayed affine formulas follow.  They are
bijective and commute with $T,R$.  On a fixed component the translation
formula is forced, and commuting with $R$ is equivalent to
$2(y-x)=0\pmod W$.  The degree difference excludes the remaining case.
These proofs include $W=1,2$ and $m=1$ even when formal affine signs
describe the same permutation.
\end{proof}

For $q>0$, take $x=x_1$.  A complete marked signature is obtained as follows.
For translation components record $j_i=(x_i-x)/d\pmod m$.
For paired components record $(e_i,j_i)$ according as
$x_i=T^{j_i}x$ or $T^{j_i}Rx$, with $e_i=0$ or $1$.
For fixed components record all $j_i=(x_i-x)/d\pmod m$ and also
\[
 k_x=(2x-a_0)/d\pmod m.
\]
The extra value records the reflection stabilizer of the root.  Equality
of $k_x,k_y$ is exactly the rooted transport condition.  In all cases
include the unmarked cover-type identifier.  For $q=0$, that identifier
alone suffices.  The API also includes the entire normalized base and
monodromy in its signature, so comparisons between different presentations
cannot silently become a claim of equivalence under a change of base.

\begin{corollary}[Exact marked-component comparison]
Within one checked presentation, two signatures are equal if and only if
there is an equivariant isomorphism respecting every ordered mark.
After preparing the cover, a signature costs $O(q)$ arithmetic operations
on binary integers.  The unique map determined by two admissible anchors
can be evaluated at any requested sheet in a constant number of such
operations.
\end{corollary}
\begin{proof}
An equivariant map preserves the root stabilizer and the relative $T,R$
coordinates.  Conversely, equal signatures satisfy the rooted transport
criterion; the resulting map takes all coordinates, hence all marks, to
their prescribed targets.  The empty-mark case is the previously proved
unmarked type classification.  No search through sheet bijections occurs.
\end{proof}

\begin{proposition}[Two marks defeat a constant type bound]
There can be $W$ inequivalent ordered two-point markings of one connected
component, even though its unmarked cover has just one type.  If the
points must be distinct, there can be $W-1$ types.
\end{proposition}
\begin{proof}
Take the annulus cover with monodromy $T(x)=x+1\pmod W$.  Any equivariant
automorphism is a translation, since the image of zero determines it.
The pairs $(0,j)$ are equivalent precisely when their second coordinates
agree.  Exclude $j=0$ for the distinct-point version.
\end{proof}
Thus polynomial bit complexity for each state and each equivalence query
does not bound the number of possible marked states by a polynomial in
$\log W$.  Nor does this example show that an unknot hierarchy must visit
all these states.  A quasi-polynomial search bound needs an additional
restriction on reachable attachment data or a way of skipping states.

\subsection{Implementation bounds, verification and measurements}

Let $L$ be the total peripheral-word length and let $B$ bound input integer
bit lengths and $\log(r+L+b+q+2)$.  The classifier uses
$O((r+L+b)B^3)$ elementary bit operations conservatively; here
$L=O(r+b)$ for the canonical schema.  It outputs
$O((r+L+b)B)$ bits.  Prepared marking queries use $O(qB^2)$ bit work
and $O(qB)$ extra storage using elementary division bounds.  The normalized
model header is shared inside a prepared index; comparing or serializing
complete signatures includes its size.  No loop depends on $W$, its
divisors, the component count, or the number of unmarked boundary lifts.

\code{CoverIndex} builds its own validated classification and returns a
defensive copy of its summary.  A serialized summary is not accepted as
a proof of a covering presentation.  Integer fields support the existing
literal-integer or signed-hexadecimal transport, including values beyond
Python's decimal conversion limit.  The standalone command supports a
cooperative deadline; cancellation propagates and never becomes a topology
answer.  Checks occur in schema construction, gcd iteration, boundary
processing and marked queries.  Individual integer operations and bulk
serialization are not hard-interruptible.

The ten maintained tests derived from the archive repeat $13846$ comparisons
against explicit lifted-spine graph search, orientation propagation, and
boundary permutation cycles.  The new independent audit exhausts all
two-generator affine actions through six sheets.  Across $364$ presentations
it performs $196712$ marked and $1061$ unmarked equivalence comparisons,
$9100$ rooted comparisons and $38746$ transported-sheet checks against
literal equivariant bijections.  Additional tests cover binary sheet counts
beyond $2^{24000}$, invalid marks, ordered repetitions, normalized input
binding, defensive copies, cancellation, and hexadecimal CLI transport.
The complete maintained suite passes $534$ tests.

\begin{center}
\begin{tabular}{@{}lrrrr@{}}
\toprule
Supplied cover & Sheets & Records & Expanded/compressed & A/A \\
\midrule
annulus & 16 & 1 & 4.06 & 1.034 \\
mobius\_reflection & 16 & 3 & 4.15 & 0.991 \\
three\_types & 16 & 3 & 3.52 & 0.990 \\
annulus & 256 & 1 & 15.65 & 1.016 \\
mobius\_reflection & 256 & 3 & 25.41 & 1.156 \\
three\_types & 256 & 3 & 17.58 & 0.988 \\
annulus & 4096 & 1 & 244.61 & 1.072 \\
mobius\_reflection & 4096 & 3 & 396.16 & 0.956 \\
three\_types & 4096 & 3 & 318.54 & 1.046 \\
annulus & 65536 & 1 & 4621.74 & 1.045 \\
mobius\_reflection & 65536 & 3 & 8290.72 & 0.989 \\
three\_types & 65536 & 3 & 7002.16 & 0.970 \\
\bottomrule
\end{tabular}
\end{center}
Ratios are median paired times for complete cover topology; sheet expansion is an independent oracle, not a competing compressed algorithm.
\begin{center}
\begin{tabular}{@{}rrrrr@{}}
\toprule
Sheet exponent & Marks & Query (ms) & JSON bytes & A/A \\
\midrule
32 & 2 & 0.010 & 100 & 1.001 \\
32 & 32 & 0.077 & 456 & 0.999 \\
32 & 512 & 1.187 & 6135 & 1.008 \\
1024 & 2 & 0.014 & 698 & 1.013 \\
1024 & 32 & 0.192 & 5534 & 1.009 \\
1024 & 512 & 3.059 & 82897 & 1.001 \\
16384 & 2 & 0.412 & 8286 & 0.998 \\
16384 & 32 & 12.330 & 70055 & 1.000 \\
16384 & 512 & 148.544 & 1058344 & 1.004 \\
\bottomrule
\end{tabular}
\end{center}
Marked queries reuse a prepared index, with sheet count $W=2^{\mathrm{exponent}}$. Preparation and serialization are outside query timing.


The topology benchmark keeps all input validation, peripheral evaluation,
orientation and genus classification inside the compressed timer.  Its expanded
arm independently constructs the lifted spine and boundary permutations.
Seven shuffled paired warm rounds retain all samples and A/A controls, with
input generation and result comparisons outside timing.  On $65536$ sheets
the observed expanded/compressed ratios range from $4622$ to $8291$; these
measure avoidance of explicit expansion, not superiority to another compressed
topology algorithm.  The topology controls range from $0.956$ to $1.156$,
so small constant-factor differences should not be overinterpreted.

The separate prepared-query experiment uses $W=2^{32},2^{1024},2^{16384}$
and $2,32,512$ ordered points spread across a paired component.  It includes
mark validation and signature construction, but excludes the shared preparation
and serialization.  At $W=2^{16384}$, $512$ marks take about $149$ milliseconds
and the JSON signature has $1058344$ bytes.  The latter illustrates that
mark output size is charged explicitly.  Query A/A controls range from
$0.998$ to $1.013$.  These finite timings support the implementation checks;
the polynomial bit-complexity claim rests on the arithmetic proof.

These measurements concern a geometric kernel with a supplied presentation.
They are not timings of unknot recognition.  To use it in a knot-exterior
hierarchy we still need a checked extraction of the base and common fibre
coordinate, an exact representation of all relevant attachments, and a
bound on the hierarchy's search and repair states.  Ordered fibre points
are only one part of that attachment interface.  The general recognizer
does not dispatch to this module or claim a sub-exponential bound from it.
